% Рисунок с доски: усечённый характеристический конус D с вершиной (x_0, t_0),
% срезанный плоскостью t = t_1. Граница состоит из трёх кусков:
% Gamma_0 (низ, t = 0), Gamma_1 (верх, t = t_1) и Gamma (боковая поверхность).
\begin{figure}[htbp]
\centering
\begin{tikzpicture}[>={Latex[length=4pt]}]
  \draw[->] (0,-0.5) -- (0,3.9) node[anchor=south east] {$t$};
  \draw[->] (-0.5,0) -- (6.3,0) node[anchor=north east] {$\R^n$};

  % усечённый конус
  \fill[pattern=north east lines, pattern color=blue!30]
    (1,0) -- (5,0) -- (4,1.5) -- (2,1.5) -- cycle;
  \draw[thick, blue!60!black] (1,0) -- (5,0);
  \draw[thick, blue!60!black] (2,1.5) -- (4,1.5);
  \draw[thick, blue!60!black] (1,0) -- (2,1.5);
  \draw[thick, blue!60!black] (5,0) -- (4,1.5);

  % срезанная верхушка
  \draw[dashed, gray] (2,1.5) -- (3,3.0) -- (4,1.5);
  \fill[gray] (3,3.0) circle (1.4pt);
  \node[gray, anchor=south, font=\footnotesize] at (3,3.08) {$(x_0,t_0)$};

  \draw[dashed, gray] (0,1.5) -- (2,1.5);
  \draw[dashed, gray] (0,3.0) -- (3,3.0);
  \node[anchor=east, font=\footnotesize] at (-0.06,1.5) {$t_1$};
  \node[anchor=east, font=\footnotesize] at (-0.06,3.0) {$t_0$};

  % нормали
  \draw[->, red!65!black] (3,0) -- (3,-0.6);
  \node[red!65!black, anchor=north, font=\footnotesize] at (3,-0.62) {$n$};
  \draw[->, red!65!black] (3,1.5) -- (3,2.1);
  \node[red!65!black, anchor=south, font=\footnotesize] at (3,2.12) {$n$};
  \draw[->, red!65!black] (4.5,0.75) -- (5.3,1.28);
  \node[red!65!black, anchor=west, font=\footnotesize] at (5.35,1.3)
    {$n = (n_1,\dots,n_n,\, n_t)$};

  \node[blue!60!black, anchor=north, font=\footnotesize] at (1.4,-0.08)
    {$\Gamma_0$};
  \node[blue!60!black, anchor=south, font=\footnotesize] at (2.4,1.54)
    {$\Gamma_1$};
  \node[blue!60!black, anchor=east, font=\footnotesize] at (1.35,0.8)
    {$\Gamma$};
  \node[blue!60!black, font=\footnotesize] at (3,0.6) {$D$};
\end{tikzpicture}
\caption{Область интегрирования в энергетическом тождестве: конус зависимости
точки $(x_0,t_0)$, срезанный на уровне $t = t_1$. На $\Gamma_0$ интеграл обнуляют
начальные данные, на $\Gamma_1$ он равен энергии, на боковой поверхности $\Gamma$
работает соотношение $n_t^2 = a^2 \sum_i n_i^2$.}
\label{fig:energy-frustum}
\end{figure}
